% Ready-to-paste journal extension for the supervisor's consensus/SHA comment.
% Required packages: booktabs, adjustbox, tikz, pgfplots.
% Required TikZ libraries: positioning, arrows.meta.

\subsection{Consensus and SHA-2 Evidence-Anchor Extension}
The supervisor-requested comparison separates two independent design layers:
the consensus protocol that orders an EVM transaction and the cryptographic
hash function that binds the evidence. SHA-256 and SHA-512 are one-way hash
functions, not encryption algorithms \cite{nistFips1804}. Ethereum's consensus
mechanism is Gasper, which combines proof-of-stake attestations, LMD-GHOST fork
choice, and Casper FFG finality \cite{ethereumGasper}. IoTeX uses Roll-DPoS,
including elected delegates, randomised committee formation, and PBFT-style
finalisation \cite{iotexRollDpos}. The existing public-testnet measurements
therefore remain correctly labelled as a Sepolia--IoTeX consensus-sensitivity
study rather than a controlled causal benchmark of consensus alone.

Figure~\ref{fig:consensus-protocols} adds Raft and QBFT without inventing
unmeasured latency. Raft is a leader-based crash-fault-tolerant replicated-log
protocol that commits through a majority quorum \cite{raftPaper,fabricRaft}.
It is not Byzantine-fault tolerant and is not directly comparable to execution
of identical Solidity bytecode on the two public EVM testnets. QBFT is an
appropriate Byzantine-fault-tolerant EVM extension target: approved Besu
validators require at least a two-thirds supermajority before inserting a block
\cite{besuQbft}. The existing ledger bytecode can be redeployed to QBFT, but the
manuscript reports only protocol properties until a real four-validator run
passes the raw-evidence gate.

\begin{figure*}[!t]
\centering
\begin{tikzpicture}[
  lane/.style={draw,rounded corners=2pt,minimum height=8mm,align=center,
               font=\scriptsize,inner sep=2pt},
  measured/.style={lane,fill=trustgreen!12,draw=trustgreen!70!black},
  extension/.style={lane,fill=trustblue!10,draw=trustblue!70!black},
  literature/.style={lane,fill=black!4,draw=black!55},
  arrow/.style={-{Latex[length=1.8mm]},thin,draw=black!65}
]
\node[measured,text width=2.55cm] (s1) at (0,0)
  {\textbf{Ethereum Sepolia}\\permissioned test validators};
\node[measured,text width=2.75cm,right=3mm of s1] (s2)
  {stake-weighted attestations\\and LMD-GHOST fork choice};
\node[measured,text width=2.55cm,right=3mm of s2] (s3)
  {Casper FFG\\economic finality};
\node[measured,text width=2.1cm,right=3mm of s3] (s4)
  {Gasper PoS\\\textbf{measured}};

\node[measured,text width=2.55cm] (i1) at (0,-1.15)
  {\textbf{IoTeX testnet}\\stake-elected delegates};
\node[measured,text width=2.75cm,right=3mm of i1] (i2)
  {randomised short-lived\\producer committee};
\node[measured,text width=2.55cm,right=3mm of i2] (i3)
  {PBFT-style committee\\finalisation};
\node[measured,text width=2.1cm,right=3mm of i3] (i4)
  {Roll-DPoS\\\textbf{measured}};

\node[literature,text width=2.55cm] (r1) at (0,-2.30)
  {\textbf{Raft}\\permissioned replicas};
\node[literature,text width=2.75cm,right=3mm of r1] (r2)
  {elected leader\\replicates ordered log};
\node[literature,text width=2.55cm,right=3mm of r2] (r3)
  {majority quorum\\deterministic commit};
\node[literature,text width=2.1cm,right=3mm of r3] (r4)
  {crash faults only\\\textbf{profiled}};

\node[extension,text width=2.55cm] (q1) at (0,-3.45)
  {\textbf{Besu QBFT}\\approved EVM validators};
\node[extension,text width=2.75cm,right=3mm of q1] (q2)
  {round-robin proposer\\and multi-round voting};
\node[extension,text width=2.55cm,right=3mm of q2] (q3)
  {$\geq 2/3$ validator\\signatures};
\node[extension,text width=2.1cm,right=3mm of q3] (q4)
  {Byzantine faults\\\textbf{not measured}};

\foreach \a/\b in {s1/s2,s2/s3,s3/s4,i1/i2,i2/i3,i3/i4,
                    r1/r2,r2/r3,r3/r4,q1/q2,q2/q3,q3/q4}
  \draw[arrow] (\a) -- (\b);
\end{tikzpicture}
\caption{Consensus-protocol comparison requested by the supervisor. Green lanes
are the two public testnets with measured first-inclusion evidence. Raft is a
crash-fault-tolerant replicated-log reference and is not an EVM public-testnet
measurement. QBFT is the selected Byzantine-fault-tolerant EVM extension target,
but no latency value is reported until an actual multi-validator Besu run passes
the same evidence gate.}
\label{fig:consensus-protocols}
\end{figure*}

\begin{table*}[!t]
\centering
\caption{SHA-256 and SHA-512 comparison used by the evidence-anchor extension.
Both are one-way SHA-2 hash functions, not encryption algorithms
\cite{nistFips1804}.}
\label{tab:sha2-semantics}
\begin{adjustbox}{width=\textwidth}
\begin{tabular}{lrrrrll}
\toprule
Algorithm & Digest & Message block & Word size & Generic collision work &
EVM execution path & Benchmark interpretation \\
& (bits) & (bits) & (bits) & (bits, approx.) & & \\
\midrule
SHA-256 & 256 & 512 & 32 & 128 &
Standard SHA-256 precompile; off-chain mode also supported &
Anchor one 32-byte full digest \\
SHA-512 & 512 & 1024 & 64 & 256 &
No standard SHA-512 precompile; compute off chain &
Anchor one 64-byte full digest \\
\bottomrule
\end{tabular}
\end{adjustbox}
\end{table*}

\begin{figure}[!t]
\centering
\begin{tikzpicture}[
  box/.style={draw,rounded corners=2pt,align=center,font=\scriptsize,
              minimum height=8mm,inner sep=2pt},
  arrow/.style={-{Latex[length=1.8mm]},thin},
  note/.style={font=\tiny,align=center,text=black!70}
]
\node[box,fill=blue!8,text width=1.45cm] (p1) {Evidence\\payload};
\node[box,fill=blue!14,text width=1.65cm,right=2.5mm of p1] (h256)
  {SHA-256\\32-byte digest};
\node[box,fill=green!12,text width=1.85cm,right=2.5mm of h256] (a256)
  {EVM anchor\\algorithm ID + digest};
\draw[arrow] (p1) -- (h256);
\draw[arrow] (h256) -- (a256);

\node[box,fill=orange!8,text width=1.45cm,below=8mm of p1] (p2)
  {Same evidence\\payload};
\node[box,fill=orange!16,text width=1.65cm,right=2.5mm of p2] (h512)
  {SHA-512\\64-byte digest};
\node[box,fill=green!12,text width=1.85cm,right=2.5mm of h512] (a512)
  {Same EVM anchor\\algorithm ID + digest};
\draw[arrow] (p2) -- (h512);
\draw[arrow] (h512) -- (a512);

\node[note,below=2mm of h512,text width=4.2cm] (boundary)
  {SHA-256 also has a standard EVM precompile.\\
   SHA-512 has no standard EVM precompile; its full digest is anchored.};
\end{tikzpicture}
\caption{Fair SHA-2 comparison boundary. Both algorithms hash identical bytes
off chain and anchor their complete digest on identical EVM bytecode. The
resulting network comparison measures digest anchoring, not encryption or CPU
hash throughput.}
\label{fig:sha2-paths}
\end{figure}


The hash benchmark computes SHA-256 and SHA-512 over identical retained payload
bytes and anchors the complete 32-byte or 64-byte digest through identical
\texttt{HashCommitmentLedger} runtime bytecode on Sepolia and IoTeX. Both
digests are computed off chain in the primary comparison, which isolates
digest-width, calldata, event, and storage effects from CPU hash throughput.
Solidity provides a SHA-256 precompile-backed built-in, whereas the standard EVM
does not provide an equivalent SHA-512 precompile \cite{solidityCrypto}. A
separate \texttt{computeSha256AndAnchor} entry point can measure the SHA-256
precompile path, but that result is not presented as a symmetric SHA-256 versus
SHA-512 computation benchmark.

\IfFileExists{derived_consensus/results_sha2_network.tex}{%
\input{derived_consensus/results_sha2_network}
\input{derived_consensus/table_sha2_network}
\input{derived_consensus/fig_sha2_anchor_gas}
}{%
No SHA-2 public-testnet number is inserted until three independent sessions
with at least 30 paired payloads each pass digest recomputation, chain-identity,
bytecode-identity, receipt, event, and negative-security validation.
}

\paragraph{Claim boundary.}
The original Fig.~6 values remain measured first-inclusion observations for
Sepolia and IoTeX only. Raft and QBFT are protocol-property comparisons unless
explicitly accompanied by measured deployment evidence. SHA-2 gas values, when
available, measure anchoring of precomputed complete digests; they do not
measure encryption, hash security, CPU throughput, consensus superiority,
economic finality, or production O-RAN performance.
